\documentclass[11pt,a4paper]{article}
\usepackage[margin=0.85in,top=1in,bottom=1in]{geometry}
\usepackage{amsmath,amssymb,amsfonts}
\usepackage{graphicx}
\usepackage{float}
\usepackage{enumitem}
\usepackage{microtype}
\usepackage{caption}
\usepackage[hidelinks]{hyperref}
\usepackage[most,breakable]{tcolorbox}
\tcbuselibrary{breakable,skins}
\usepackage[utf8]{inputenc}
\usepackage[T1]{fontenc}
\usepackage{lmodern}
\captionsetup{font=small,labelfont=bf,skip=4pt}
\setlist{leftmargin=1.5em,topsep=2pt}

%% Breakable card — prevents content cutoff across pages
\newtcolorbox{solcard}[3]{
  breakable, enhanced,
  colback=white, colframe=gray!50, arc=2pt,
  title={\textbf{#1}\hfill\textsf{\small #3\,pts}},
  fonttitle=\normalsize\sffamily\bfseries,
  before upper={\small\textit{#2}\par\smallskip\hrule height 0.4pt\smallskip},
  top=4pt, bottom=6pt, left=7pt, right=7pt,
  parbox=false,
}

\title{\textbf{Domino effect --- Solution}}
\author{\normalsize Y25 (2025) — English translation}
\date{}
\begin{document}\maketitle\vspace{-0.5em}
\begin{solcard}{A1}{Measure the sequence $t_n$ for $n\in\{1,3,5,7,9,11\}$ at $s=2~cm$. Plot the dependence $t_n(n)$.}{0.70}
\par\smallskip\noindent\textbf{Answer:} $n$ 1 3 5 7 9 11 $t_n,~\mathrm{ms}$ 0 160 222 272 314 356
\begin{figure}[H]
  \centering
  \includegraphics[width=0.72\textwidth]{E1_S_EN_assets/fig1.jpg}
\end{figure}
\end{solcard}

\begin{solcard}{A2}{Express the angles $\beta$ and $\varphi$ in terms of the geometric parameters $s$, $h$, $\alpha$.}{0.10}
\par\smallskip\noindent\textbf{Answer:} \[\beta = \operatorname{arctg}\left(\dfrac{d}{h}\right)\]\[\varphi = \arcsin\left(\dfrac{s-d}{h}\right) - \operatorname{arctg}\left(\dfrac{d}{h}\right)\]
\end{solcard}

\begin{solcard}{A3}{Express $\Omega_0$ in terms of $g,h,k, \alpha, s$ in a steady wave of falling dominoes. Please note that at the moment of impact, instantaneous reaction forces arise at all points of contact of the knuckles with each other and with the surface.}{2.50}

The condition of the kinematic connection immediately after the impact: \[(\vec{v}_M\cdot\vec{e}_x)=(\vec{v}_N\cdot\vec{e}_x)\]Substituting $\vec{v}_M = \left[\vec{\Omega}_a\times \vec{r}_M\right]$, $\vec{v}_N = \left[\vec{\Omega}_0\times \vec{r}_N\right]$, we get: \[\Omega_0 y = \Omega_a y \Rightarrow \Omega_0=\Omega_a\]


From an energy point of view, the system at times $t_n$ and $t_{n+1}$ differs in that one domino at rest moves to the beginning of the row. Based on this, we can write down the change in potential energy: $\Delta \Pi = \dfrac{mgh}{2}\left(1 - \sqrt{1+\alpha^2}\sin\left(\operatorname{arctg}\alpha+\arcsin \dfrac{d}{s}\right)\right)$


Let's write down the law of conservation of energy, taking into account the kinematic connection: \[\dfrac{I\Omega_b^2}{2} =2\cdot \dfrac{I\Omega_0^2}{2}+\dfrac{mgh}{2}\left(1 - \sqrt{1+\alpha^2}\sin\left(\operatorname{arctg}\alpha+\arcsin \dfrac{d}{s}\right)\right),\]where $I= mh^2(1+\alpha^2)/3$ is the moment of inertia of the domino relative to the axis of rotation. Substituting $\Omega_b=k\Omega_0$, we get: \[(k^2-2)\Omega_0^2=\dfrac{3g\left(1 - \sqrt{1+\alpha^2}\sin\left(\operatorname{arctg}\alpha+\arcsin \dfrac{d}{s}\right)\right)}{h(1+\alpha^2)}\]

\par\smallskip\noindent\textbf{Answer:} \[\Omega_0=\sqrt{\dfrac{3g\left(1 - \sqrt{1+\alpha^2}\sin\left(\operatorname{arctg}\alpha+\arcsin \dfrac{\alpha h}{s}\right)\right)}{h(1+\alpha^2)(k^2-2)}}\]
\begin{figure}[H]
  \centering
  \includegraphics[width=0.72\textwidth]{E1_S_EN_assets/fig2.jpg}
\end{figure}
\begin{figure}[H]
  \centering
  \includegraphics[width=0.72\textwidth]{E1_S_EN_assets/fig3.jpg}
\end{figure}
\end{solcard}

\begin{solcard}{B1}{Take the necessary measurements and obtain $c$ values for at least 10 different $s/h$. Indicate the error of the found $c$.}{4.30}

The speed of propagation of wave $c$ is equal to: \[c=\dfrac{s}{\mathrm{d}t_n/\mathrm{d}n}\] When calculating $\mathrm{d}t_n/\mathrm{d}n$ such $n$ should be used for which the dependence $t_n(n)$ is linear. But at the same time, the number of points used when calculating $\mathrm{d}t_n/\mathrm{d}n$ must be at least 4.


\[c=\dfrac{s}{\mathrm{d}t_n/\mathrm{d}n}\]


When calculating $\mathrm{d}t_n/\mathrm{d}n$, those $n$ should be used for which the dependence $t_n(n)$ is linear. But at the same time, the number of points used when calculating $\mathrm{d}t_n/\mathrm{d}n$ must be at least 4.


$s,~ cm$ $\mathrm{d}t_n/\mathrm{d}n, ~\mathrm{ms}$ $c,~m/s$ 1.0 $n$ 7 9 11 12 14 16 9.55 1.047 $t_n,~\mathrm{ms}$ 0 30 57.6 76.8 96 115 1.2 $n$ 4 6 8 10 12 14 14.30 0.839 $t_n,~\mathrm{ms}$ 0 47 91 112 149 174 1.7 $n$ 10 11 12 13 14 15 15.80 1.076 $t_n,~\mathrm{ms}$ 0 14 33.6 43.6 63.2 79.2 2.0 $n$ 8 9 10 11 12 13 19.00 1.053 $t_n,~\mathrm{ms}$ 0 27.2 47.2 62.4 85.6 103 2.3 $n$ 6 7 8 9 11 12 23.90 0.962 $t_n,~\mathrm{ms}$ 0 35.6 62.4 80.4 128 158 2.5 $n$ 3 5 6 7 8 9 27.10 0.923 $t_n,~\mathrm{ms}$ 0 62 88 115.6 142.8 170 2.7 $n$ 4 6 7 8 9 10 31.30 0.863 $t_n,~\mathrm{ms}$ 0 63 91 125 152 189 3.0 $n$ 2 3 4 5 6 7 50.60 0.652 $t_n,~\mathrm{ms}$ 0 58 116 170 210 252 3.3 $n$ 3 4 5 6 7 9 41.50 0.795 $t_n,~\mathrm{ms}$ 0 61 106 146 193 268 3.5 $n$ 3 4 5 6 7 8 61.90 0.565 $t_n,~\mathrm{ms}$ 0 58 122 188 242 310


Two main factors contribute to the error in finding $c$: the error in determining the time $t_n$, associated with the imperfection of the detector: when a domino falls, it does not instantly cover the light incident on the photodiode. This appears as non-ideal steps with a transition width of $\approx 4~\mathrm{ms}$. Considering that the duration of the entire process of wave passage is $\approx 200-300~\mathrm{ms}$, we have an error of $\approx 5 \%$. Another factor of error is associated with the fact that each time we place the dominoes slightly differently (which are also different from each other). As a result, this leads to the fact that the measured speed $c$ may differ by $\approx 5\%$ when repeated.


This appears as non-ideal steps with a transition width of $\approx 4~\mathrm{ms}$. Considering that the duration of the entire process of wave passage is $\approx 200-300~\mathrm{ms}$, we have an error of $\approx 5 \%$.


Another factor of error is associated with the fact that each time we place the dominoes slightly differently (which are also different from each other). As a result, this leads to the fact that the measured speed $c$ during repetition may differ by $\approx 5\%$.

\begin{figure}[H]
  \centering
  \includegraphics[width=0.72\textwidth]{E1_S_EN_assets/fig4.jpg}
\end{figure}
\end{solcard}

\begin{solcard}{B2}{Plot a graph of $k$ versus $s/h$.}{2.40}

Let us denote by $\tau = \mathrm{d}t/\mathrm{d}n$ \[\dfrac{s\Omega_0}{2(\beta+\varphi)}\left( 1+\dfrac{1}{k}\right) = \dfrac{s}{\tau}\] Let us substitute the expression for $\Omega_0$ from point A3: \[\dfrac{1}{2(\beta+\varphi)}\left( 1+\dfrac{1}{k}\right)\cdot \sqrt{\dfrac{3g\left(1 - \sqrt{1+\alpha^2}\sin\left(\operatorname{arctg}\alpha+\arcsin \dfrac{d}{s}\right)\right)}{h(1+\alpha^2)(k^2-2)}} = \dfrac{1}{\tau}\] \[\dfrac{\tau}{2(\beta+\varphi)}\cdot\sqrt{\dfrac{3g\left(1 - \sqrt{1+\alpha^2}\sin\left(\operatorname{arctg}\alpha+\arcsin \dfrac{d}{s}\right)\right)}{h(1+\alpha^2)}}\cdot \dfrac{1+\dfrac{1}{k}}{\sqrt{k^2-2}}=1\] We denote the factor on the left by $u$: \[u = \dfrac{\tau}{2(\beta+\varphi)}\cdot\sqrt{\dfrac{3g\left(1 - \sqrt{1+\alpha^2}\sin\left(\operatorname{arctg}\alpha+\arcsin \dfrac{d}{s}\right)\right)}{h(1+\alpha^2)}}\] Then the equation will look like this: \[u(s/h) \cdot \dfrac{1+\dfrac{1}{k}}{\sqrt{k^2-2}}=1,\] which can be solved by iteration. First, for each of the measured $s$ we calculate $u$, then using the iteration method we find the value $k$: \[k = \sqrt{2+u^2\left(1+\dfrac{1}{k}\right)^2}\]


\[\dfrac{\tau}{2(\beta+\varphi)}\cdot\sqrt{\dfrac{3g\left(1 - \sqrt{1+\alpha^2}\sin\left(\operatorname{arctg}\alpha+\arcsin \dfrac{d}{s}\right)\right)}{h(1+\alpha^2)}}\cdot \dfrac{1+\dfrac{1}{k}}{\sqrt{k^2-2}}=1\]


We denote the factor on the left as $u$: \[u = \dfrac{\tau}{2(\beta+\varphi)}\cdot\sqrt{\dfrac{3g\left(1 - \sqrt{1+\alpha^2}\sin\left(\operatorname{arctg}\alpha+\arcsin \dfrac{d}{s}\right)\right)}{h(1+\alpha^2)}}\] Then the equation will look like this: \[u(s/h) \cdot \dfrac{1+\dfrac{1}{k}}{\sqrt{k^2-2}}=1,\] which can be solved by iteration. First, for each of the measured $s$ we calculate $u$, then using the iteration method we find the value $k$:


\[k = \sqrt{2+u^2\left(1+\dfrac{1}{k}\right)^2}\]


$s,~ cm$ $s/h$ $\tau,~ \mathrm{ms}$ $\varphi, \mathrm{rad}$ $u$ $k$ 1.0 0.25 9.55 -0.024 0.659 1.753 1.2 0.3 10.8 0.027 0.576 1.686 1.7 0.425 15.80 0.155 0.538 1.657 2.0 0.5 19.00 0.234 0.530 1.651 2.3 0.575 23.90 0.316 0.562 1.675 2.5 0.625 27.10 0.373 0.575 1.685 2.7 0.675 31.30 0.431 0.603 1.707 3.0 0.75 38.80 0.523 0.653 1.748 3.3 0.825 50.60 0.621 0.749 1.828 3.5 0.875 61.90 0.693 0.843 1.910

\begin{figure}[H]
  \centering
  \includegraphics[width=0.72\textwidth]{E1_S_EN_assets/fig5.jpg}
\end{figure}
\end{solcard}

\end{document}